\documentclass[twocolumn,10pt,notitlepage]{article}
\usepackage{bm}
\usepackage{breqn}
\usepackage{diagbox}
\usepackage{float}
\usepackage{xcolor}


\title{A Fast and Space-Efficient Algorithm for Validating Pointer Array Formats}
\date{2022-07-27}
\author{Malcolm Roberts, \texttt{malcolm.roberts@amd.com}}

\begin{document}

\maketitle

\section{Abstract}
Pointer arithmetic lies at the heart of high-performance numerical
algorithms on CPUs and GPUs.  A $d$-dimensional array is specified by
providing a length $\bm{\ell}=\left(\ell_0, \dots, \ell_{d-1}\right)$,
and strides $\bm{s}=\left(s_0, \dots, s_{d-1}\right)$. Data at
multi-index $\bm{a}=\left(a_0,\dots,a_{d-1} \right)$, with
$0\leq{}a_i<\ell_i$, is stored at pointer offset $\bm{a}\cdot\bm{s}$.
A data format is considered invalid if two different indices point to
the same memory location.  That is, we require that the map from
multi-index to memory address be injective.  This can be verified
using a hash table, but the complexity is
$\mathcal{O}\left(\ell^d\right)$ in time and space, which can be
prohibitively expensive.  In this paper, we present an algorithm which
reduces the time complexity to $\mathcal{O}\left(\ell^{d-1}\right)$
and the space complexity to $\mathcal{O}\left(\ell^{d-2}\right)$.  We
demonstrate how this can be used in ROCm math libraries to produce
tests parameters and validate user input data.


\section{Introduction}

Data management is at the heart of high-performance computing.  On the
GPU, kernels access data via pointers passed as arguments.  The
pointer refers to an the address of a buffer in some flat memory
space; structure is aded to this data by passing additional
parameters.  From the programmer's perspective, all data buffers are
one-dimensional.  In the case of single- or multi-dimensional arrays,
the data pointer is accompanied by length and stride information.

\subsection{Scalar and 1D}
Let $\texttt{p}$ be a pointer to the address of a data buffer in
memory. If we wish to interpret the data in this buffer as a
one-dimensional array, we associate a integer length $\ell$ and an
integer stride $s$ to the pointer $\texttt{p}$.

If $\ell=0$, the data is scalar, and we can access the data at
$\texttt{p[0]}$ without any restrictions on the stride $s$.  

If $\ell>0$ then we wish to access the data at index $a$ with
$a\in\left[0,\dots,\ell-1\right)$.  The data corresponding to index
  $a$ is stored in memory at $\texttt{p[}a \times{} s\texttt{]}$.  The
  data format for a one-dimensional array is given by pair
  $\left(\ell, s\right)$, and we consider this data format to be valid
  if and only if the map
\begin{dmath}
  \label{onedmap}
  a \rightarrow a \times s
\end{dmath}
is injective for $a\in\left[0,\dots,\ell-1\right)$.  That is, each
  unique index $a$ must refer to a distinct memory address.  If this
  map is not injective, then two different indices will address the
  same memory location, resulting in data loss.
  
The one-dimensional is quite simple: we are guaranteed that the map in
equation~\eqref{onedmap} is injective if and only if $s$ is non-zero.
However, this becomes more complicated in multiple dimensions.


\subsection{Higher dimensions}
For $d$-dimensional arrays, the length and stride are given by
length-$d$ vectors, $\bm{\ell}=\left(\ell_0, \dots, \ell_{d-1}\right)$
and $\bm{s}=\left(s_0, \dots, s_{d-1}\right)$.  We take all of these
values to be strictly positive; lengths are always non-negative, and
if any $\ell_i=0$ then we can kick out that index and consider the
$(d-1)$-dimensional case.  Moreover, we insist that $s_i>0$; while
negative strides can be useful, we will treat them separately.  The
multi-index $\bm{a}=\left(a_0,\dots,a_{d-1}\right)$, with
$0\leq{}a_i<\ell_i$, references data at
$\texttt{p[}\bm{a}\cdot{}\bm{s}\texttt{]}$.

Consider a 2D array with strides $\bm{s}=(3,5)$.  For length
$\bm{\ell}=\left(5,4\right)$, the address offsets are given in
Table~\ref{tab:2dok}, and it can be easily seen that all memory
addresses are unique.  Increasing the length to
$\bm{\ell}=\left(6,4\right)$, we have a memory address collision, as
shown in Table~\ref{tab:2dbad}.  In particular,
$(5,0)\rightarrow{}5\times{}3 + 0 \times{}5=15$ and
$(0,3)\rightarrow{}0\times{}3 + 3 \times{}5=15$ coincide. At first
glance, it's not immediately obvious that $\left((5,4),(3,5)\right)$
would lead to a valid array format, and that
$\left((6,4),(3,5)\right)$ would not.  Tellingly, the collision occurs
at the edge of the table, which we will demonstrate to be the case in
general.

\begin{table}
  \centering
  \begin{tabular}{|l|l|l|l|l|l|l|}
    \hline
    \backslashbox{$a_1$}{$a_0$}& 0  & 1  & 2  & 3  & 4 \\
    \hline 0 & 0  & 3  & 6  & 9  & 12 \\
    \hline 1 & 5  & 8  & 11 & 14 & 17 \\
    \hline 2 & 10 & 13 & 16 & 19 & 22 \\
    \hline 3 & 15 & 18 & 21 & 24 & 27 \\
    \hline 
  \end{tabular}
  \caption{\label{tab:2dok}2D array with
    $(\bm{\ell},\bm{s})=\left((5,4),(3,5)\right)$.}
\end{table}
\begin{table}
  \centering
  \begin{tabular}{|l|l|l|l|l|l|l|l|}
    \hline
    \backslashbox{$a_1$}{$a_0$}& 0  & 1  & 2  & 3  & 4 & 5\\
    \hline 0 & 0  & 3  & 6  & 9  & 12 & {\color{red}15} \\
    \hline 1 & 5  & 8  & 11 & 14 & 17 & 20 \\
    \hline 2 & 10 & 13 & 16 & 19 & 22 & 25 \\
    \hline 3 & {\color{red}15} & 18 & 21 & 24 & 27 & 30 \\
    \hline 
  \end{tabular}
  \caption{\label{tab:2dbad}2D array with
    $(\bm{\ell},\bm{s})=\left((6,4),(3,5)\right)$.}
\end{table}

\section{Validation algorithms}

Given a pair $(\bm{\ell},\bm{s})$, the direct computation of all
allowble indices has involves computing $N\doteq{}\prod_i \ell_i$ dot
producs, which requires $\mathcal{O}(dN)$ multiplications.  The
required storage is of size $N$.  A direct search for duplicates
requires $\mathcal{O}(N^2)$ operations; this can be improved to
$\mathcal{O}(N)$ by using a hash table.  The cost of computing and
storing all of the memory offsets is prohibitive: $N$ may be of the
order of a few billion, and the validation of the data format may be
as expensive as the the actual operation to be performed on the data!

Approaching the validation problem from a number-theoretic perspective
allows us to avoid most of this cost.

\subsection{2D coprime case}



\end{document}
